\documentclass{beamer}
\mode<presentation>
\usepackage{amsmath}
\usepackage{amssymb}
\usepackage{adjustbox}
\usepackage{subcaption}
\usepackage{enumitem}
\usepackage{multicol}
\usepackage{mathtools}
\usepackage{listings}
\usepackage{float}
\usepackage{graphicx}
\usepackage{url}
\def\UrlBreaks{\do\/\do-}

\usetheme{Boadilla}
\usecolortheme{lily}

\setbeamertemplate{footline}{
  \leavevmode%
  \hbox{%
  \begin{beamercolorbox}[wd=\paperwidth,ht=2.25ex,dp=1ex,right]{author in head/foot}%
    \insertframenumber{} / \inserttotalframenumber\hspace*{2ex}
  \end{beamercolorbox}}%
  \vskip0pt%
}
\setbeamertemplate{navigation symbols}{}

% ----- Macros from your template -----
\providecommand{\nCr}[2]{\,^{#1}C_{#2}}
\providecommand{\nPr}[2]{\,^{#1}P_{#2}}
\providecommand{\mbf}{\mathbf}
\providecommand{\pr}[1]{\ensuremath{\Pr\left(#1\right)}}
\providecommand{\qfunc}[1]{\ensuremath{Q\left(#1\right)}}
\providecommand{\sbrak}[1]{\ensuremath{{}\left[#1\right]}}
\providecommand{\lsbrak}[1]{\ensuremath{{}\left[#1\right.}}
\providecommand{\rsbrak}[1]{\ensuremath{\left.#1\right]}}
\providecommand{\brak}[1]{\ensuremath{\left(#1\right)}}
\providecommand{\lbrak}[1]{\ensuremath{\left(#1\right.}}
\providecommand{\rbrak}[1]{\ensuremath{\left.#1\right)}}
\providecommand{\cbrak}[1]{\ensuremath{\left\{#1\right\}}}
\providecommand{\lcbrak}[1]{\ensuremath{\left\{#1\right.}}
\providecommand{\rcbrak}[1]{\ensuremath{\left.#1\right\}}}
\newtheorem{rem}{Remark}
\newcommand{\sgn}{\mathop{\mathrm{sgn}}}
\providecommand{\abs}[1]{\left\vert#1\right\vert}
\providecommand{\res}[1]{\Res\displaylimits_{#1}}
\providecommand{\norm}[1]{\lVert#1\rVert}
\providecommand{\mtx}[1]{\mathbf{#1}}
\providecommand{\mean}[1]{E\left[ #1 \right]}
\providecommand{\fourier}{\overset{\mathcal{F}}{ \rightleftharpoons}}
\providecommand{\system}{\overset{\mathcal{H}}{ \longleftrightarrow}}
\providecommand{\dec}[2]{\ensuremath{\overset{#1}{\underset{#2}{\gtrless}}}}
\newcommand{\myvec}[1]{\ensuremath{\begin{pmatrix}#1\end{pmatrix}}}
\let\vec\mathbf

\lstset{language=C,frame=single,breaklines=true,columns=fullflexible}
\numberwithin{equation}{section}

\title{Problem 12.526 — Jordan Canonical Form \& GM}
\author{Aryansingh Sonaye \\ AI25BTECH11032 \\ EE1030 - Matrix Theory}
\date{\today}

\begin{document}

\begin{frame}
\titlepage
\end{frame}

\section{Problem}
\begin{frame}
\frametitle{Problem 12.526}
Let $A$ be a $7\times7$ matrix such that 
\begin{align}
2A^2 - A^4 = I_7
\end{align}
where $I_7$ is the identity matrix. If $A$ has two distinct eigenvalues and each eigenvalue has geometric multiplicity $3$, then the total number of nonzero entries in the Jordan canonical form of $A$ equals:
\begin{align*}
\text{(A) } 6 \qquad \text{(B) } 7 \qquad \text{(C) } 8 \qquad \text{(D) } 9
\end{align*}
\end{frame}

\section{Given Data}
\begin{frame}
\frametitle{Given Data}
\begin{table}[H]\centering
\renewcommand{\arraystretch}{1.4}
\begin{tabular}{|c|c|c|}
\hline
\textbf{Symbol} & \textbf{Meaning} & \textbf{Value} \\
\hline
$A$ & Square matrix & $7\times7$ \\
\hline
Equation & Polynomial relation & $2A^2 - A^4 = I_7$ \\
\hline
Eigenvalues & Distinct eigenvalues of $A$ & Two distinct ($\lambda=\pm1$) \\
\hline
Geometric multiplicity (GM) & $\dim(\ker(A-\lambda I))$ & $3$ for each eigenvalue \\
\hline
Total GM & Sum over both eigenvalues & $3+3=6$ \\
\hline
\end{tabular}
\end{table}
\end{frame}

\section{Theory}
\subsection{Geometric Multiplicity}
\begin{frame}
\frametitle{1. Geometric Multiplicity}
For a square matrix $A$ and eigenvalue $\lambda$,
\begin{align}
\text{GM}(\lambda) = \dim\!\big(\ker(A-\lambda I)\big).
\end{align}
It represents the number of linearly independent eigenvectors for $\lambda$.

If the \textbf{algebraic multiplicity (AM)} equals GM, then $A$ is diagonalizable; otherwise $A$ is \textbf{defective} and has Jordan blocks larger than $1\times1$.
\end{frame}

\subsection{Jordan Canonical Form}
\begin{frame}
\frametitle{2. Jordan Canonical Form}
Every square matrix $A$ is similar to a matrix $J$ such that
\begin{align}
A = PJP^{-1}
\end{align}
where $J$ is the \textbf{Jordan canonical form} and $P$ has columns that are eigenvectors and generalized eigenvectors.

A Jordan block of size $k$ for eigenvalue $\lambda$ is:
\begin{align}
J_k(\lambda) =
\myvec{
\lambda & 1 & 0 & \cdots & 0\\
0 & \lambda & 1 & \cdots & 0\\
\vdots & & \ddots & \ddots & \vdots\\
0 & \cdots & 0 & \lambda & 1\\
0 & \cdots & \cdots & 0 & \lambda
}.
\end{align}
The $1$ above the diagonal links a generalized eigenvector to a true eigenvector, e.g.
\begin{align}
\myvec{\lambda & 1\\0 & \lambda},
\end{align}
with $(A-\lambda I)\vec{v_1}=0$ and $(A-\lambda I)\vec{v_2}=\vec{v_1}$.
\end{frame}

\subsection{Null-space Ladder}
\begin{frame}
\frametitle{3. Determining Block Sizes via Null Spaces}
Define
\begin{align}
N_k(\lambda)=\ker\!\big((A-\lambda I)^k\big).
\end{align}
These satisfy
\begin{align}
N_1(\lambda)\subseteq N_2(\lambda)\subseteq N_3(\lambda)\subseteq\cdots
\end{align}
and their dimensions reveal Jordan structure:

\begin{table}[H]\centering
\renewcommand{\arraystretch}{1.3}
\begin{tabular}{|c|c|}
\hline
\textbf{Quantity} & \textbf{Meaning} \\ \hline
$\dim N_1(\lambda)$ & Number of Jordan blocks (= GM) \\ \hline
$\dim N_2-\dim N_1$ & Number of blocks of size $\ge2$ \\ \hline
$\dim N_3-\dim N_2$ & Number of blocks of size $\ge3$ \\ \hline
\end{tabular}
\end{table}
Thus, computing $\ker\big((A-\lambda I)^k\big)$ for successive $k$ tells how many blocks of each size exist.
\end{frame}

\section{Solution}
\subsection{Polynomial \& Eigenvalues}
\begin{frame}
\frametitle{Step 1: Polynomial \& Step 2: Eigenvalues}
\begin{align}
2A^2 - A^4 = I_7
&\Rightarrow A^4 - 2A^2 + I = 0\\
&\Rightarrow (A^2 - I)^2 = 0
\end{align}
Let $\lambda$ be an eigenvalue of $A$.
\begin{align}
(\lambda^2 - 1)^2 = 0 \Rightarrow \lambda^2 = 1 \Rightarrow \lambda = \pm 1.
\end{align}
So $A$ has two distinct eigenvalues: $+1$ and $-1$.
\end{frame}

\subsection{GM data \texorpdfstring{$\Rightarrow$}{=>} structure}
\begin{frame}
\frametitle{Step 3: GM Data $\Rightarrow$ Jordan Structure}
\begin{align}
\text{GM}(+1)=3,\qquad \text{GM}(-1)=3.
\end{align}
Total eigenvectors $=6$. Since $A$ is $7\times7$, one dimension is missing — corresponding to one generalized eigenvector.\\[4pt]
$\Rightarrow$ There is \textbf{one} Jordan block of size $2$ and \textbf{five} blocks of size $1$.
\end{frame}

\subsection{Jordan matrix}
\begin{frame}
\frametitle{Step 4: Jordan Form}
(Assume the $2\times2$ block belongs to $\lambda=+1$; the count is unaffected.)
\begin{align}
J =
\myvec{
1 & 1 & 0 & 0 & 0 & 0 & 0\\
0 & 1 & 0 & 0 & 0 & 0 & 0\\
0 & 0 & 1 & 0 & 0 & 0 & 0\\
0 & 0 & 0 & 1 & 0 & 0 & 0\\
0 & 0 & 0 & 0 & -1 & 0 & 0\\
0 & 0 & 0 & 0 & 0 & -1 & 0\\
0 & 0 & 0 & 0 & 0 & 0 & -1
}.
\end{align}
\end{frame}

\subsection{Counting nonzeros}
\begin{frame}
\frametitle{Step 5: Count Nonzero Entries}
Each $1\times1$ block contributes one nonzero.  
The $2\times2$ block contributes three (two diagonal entries and one superdiagonal):
\begin{align}
\text{Total nonzeros}=5\times1+3=8.
\end{align}
\begin{block}{\centering Final Answer}
\Large \(\boxed{8}\)
\end{block}
\end{frame}

\section{Conclusion}
\begin{frame}
\frametitle{Conclusion}
The null spaces $\ker\!\big((A-\lambda I)^k\big)$ reveal the Jordan block sizes: each jump in $\dim N_k$ counts blocks of size $\ge k$.  
Here, only one eigenvector is missing from full diagonalizability $\Rightarrow$ exactly one $2\times2$ block.

\begin{align}
\boxed{\text{Total nonzero entries in the Jordan form of }A=8.}
\end{align}
\end{frame}

\end{document}

